% Check criterion, deliverable, and the Decision Log question of Day 7.
% The criterion and the deliverable come from the Day 7 section of the
% syllabus and from Labs/day07/README.md.
\begin{frame}{Check criterion}
  \small
  The instructor checks this at the end of the lab:
  \begin{itemize}
    \item \cmark\ \code{python pi/classify\_image.py images/Cat03.jpg} prints
          \code{Task B1: complete} and a cat class as the first class.
    \item \cmark\ The notebook prints \code{Tasks complete: 1, 2, 3} in its
          last cell.
    \item \cmark\ \code{python pi/bench.py} prints \code{Task D1: complete}.
    \item \cmark\ The latency table of the report has two runtimes or more,
          two thread counts or more, and two precisions or more.
    \item \cmark\ The report has the five answers of the Netron questions.
    \item \cmark\ The report has a prediction and a result for each question
          of Part D.
    \item \cmark\ The Decision Log gives numbers and names one trade-off.
  \end{itemize}
  \vskip 0.4em
  \textbf{Hand in:} the file \code{report.md}.
\end{frame}

\begin{frame}{Decision Log}
  Write about 100 words. State one design decision, give your measured
  numbers, and name the trade-off.
  \vskip 0.4em
  \begin{exerciseblock}
    A product with a Raspberry Pi 5 and a camera must classify 10 images in
    each second. The same program also reads the camera and sends each
    result over the network.
    \begin{itemize}
      \item Which runtime, which precision, and how many threads do you
            select for the model?
    \end{itemize}
  \end{exerciseblock}
  \vskip 0.2em
  A good Decision Log:
  \begin{itemize}
    \item gives the latency of your selection and of one other setting from
          your table,
    \item compares the latency with the budget of 100 ms for each image,
    \item names what you lose with your decision. Example: the \code{int8}
          file is faster, and it loses about 2 points of accuracy.
  \end{itemize}
\end{frame}
